\documentclass[11pt]{article}
\usepackage[margin=1in]{geometry}
\usepackage{amsmath,amssymb,amsthm}
\usepackage[colorlinks=true,urlcolor=blue,linkcolor=blue,citecolor=blue]{hyperref}
\setlength{\emergencystretch}{1em}
\newtheorem{theorem}{Theorem}
\newtheorem{lemma}{Lemma}
\title{Positive coefficients for the elliptic-integral\\upper remainder and its companion functions}
\author{Principia Math in Progress}
\date{12 September 2026}
\begin{document}
\maketitle

\paragraph{Task 1: Resolve MathDB 350004.}
We complete Conjecture 5.1 of Qiu--Bao--Ma--Jiang \cite{QBMJ}.
The principal new step is positivity of every coefficient of the
upper remainder \(h_{11}\). For the quotient \(G\), we explicitly use
the recently published theorem of Yang--Zhao \cite{YZ}; the resulting
proof does not claim that already established part as new.

Put \(x=r^2\), and write
\[
 F_0(x)={}_2F_1(\tfrac12,\tfrac12;1;x)
       =\frac2\pi\mathcal K(\sqrt x),\qquad
 F_1(x)={}_2F_1(\tfrac12,1;\tfrac32;x)
       =\frac{\operatorname{arctanh}\sqrt x}{\sqrt x}.
\]
Let \(A=\log F_0\), \(L=\log F_1\), and \(Q=A/L\), with the
removable value \(Q(0)=3/4\). The six functions in the conjecture are
\begin{align*}
 G(r)&=Q(x),&
 f(r)&=\frac{Q(x)-3/4}{x},&
 f_7(r)&=\frac{L(x)}x,\\
 h_{11}(r)&=\frac{H(x)}{x^2},&
 H(x)&=\frac{3+x}{4}L(x)-A(x),&
 h_{13}(r)&=\frac{A(x)}x,\\
 h_{12}(r)&=\frac{A(x)-(3/4+x/320)L(x)}{x^3}.
\end{align*}

\begin{theorem}\label{thm:main}
Each of \(f,G,f_7,h_{11},h_{12},h_{13}\) has a convergent expansion
\(\sum_{n\geq0}c_nr^{2n}\) for \(0\leq r<1\), with \(c_n>0\)
for every \(n\). In particular all six functions are strictly
absolutely monotone on \((0,1)\): every derivative, including order
zero, is positive there. The function \(f\) is strictly convex, and
\(h_{11},h_{12}\) are strictly increasing and strictly convex, with ranges
\[
 h_{11}((0,1))=(79/960,\log(\pi/2)),\qquad
 h_{12}((0,1))=(517/604800,\infty).
\]
\end{theorem}

As the displayed source expansions make clear, positive
coefficients refers to the even-power series. Odd powers have
zero coefficients, since all six functions are even.

\paragraph{A coefficient-sign lemma.}
We give the elementary reciprocal argument underlying the proof,
so no unproved coefficient-positivity criterion is needed.

\begin{lemma}\label{lem:sign}
Suppose \(B(x)=\sum_{n\geq0}b_nx^n\), \(b_0=1\), \(b_n>0\),
has radius of convergence at least one, and \(b_n/b_{n-1}\) is
nondecreasing. Then
\[
 \frac1{B(x)}=1-\sum_{n\geq1}c_nx^n,\qquad 0\leq c_n\leq b_n.
\]
If the ratios are strictly increasing, every \(c_n\) is positive.
If \(U(x)=\sum_{n\geq0}u_nx^n\) has positive coefficients and
\(u_n/b_n\) is nonincreasing, every nonconstant coefficient of \(U/B\)
is nonpositive. It is strictly negative when both ratio
conditions are strict.
\end{lemma}
\begin{proof}
Coefficient comparison determines
\[
 c_n=b_n-\sum_{j=1}^{n-1}c_jb_{n-j}.
\]
Inductively, \(c_1=b_1>0\), and
\[
 \sum_{j=1}^{n-1}c_jb_{n-j}
 \leq\frac{b_n}{b_{n-1}}
      \sum_{j=1}^{n-1}c_jb_{n-1-j}
 =b_n.
\]
Thus \(0\leq c_n\leq b_n\). For strictly increasing ratios the
inequality is strict when \(n\geq2\), already in the \(j=1\) term.
The bound \(c_n\leq b_n\) ensures convergence on the unit disk.
For \(n\geq1\), set \(\rho_n=u_n/b_n\). Then
\[
 [x^n]\frac{U}{B}
 =u_n-\sum_{j=1}^nc_ju_{n-j}
 =\sum_{j=1}^nc_jb_{n-j}(\rho_n-\rho_{n-j})\leq0.
\]
Strictness follows from the \(j=n\) term.
\end{proof}

\paragraph{Three quotients with negative coefficients.}
Define
\[
 G_0(x)={}_2F_1(\tfrac12,\tfrac12;2;x),\qquad
 G_1(x)={}_2F_1(\tfrac12,1;\tfrac52;x).
\]
Direct series comparison gives
\begin{equation}\label{eq:derivatives}
 F_0'=\frac{G_0}{4(1-x)},\qquad
 F_1'=\frac{G_1}{3(1-x)}.
\end{equation}
Write \(a_n=( (1/2)_n/n!)^2\). The four coefficient sequences are
\[
 [x^n]F_0=a_n,\quad [x^n]F_1=\frac1{2n+1},\quad
 [x^n]G_0=\frac{a_n}{n+1},\quad
 [x^n]G_1=\frac3{(2n+1)(2n+3)}.
\]
The denominator sequences of \(G_0/F_0\), \(G_1/F_1\), and \(G_0/G_1\)
are strictly log-convex: their consecutive ratios are respectively
\[
 \left(\frac{n-1/2}{n}\right)^2,\qquad
 \frac{2n-1}{2n+1},\qquad
 \frac{2n-1}{2n+3}\qquad(n\geq1),
\]
all strictly increasing.
For the first two quotients, the ratios of numerator to denominator
coefficients are \(1/(n+1)\) and \(3/(2n+3)\), both strictly decreasing.
For the third, if \(\rho_n=[x^n]G_0/[x^n]G_1\), then
\[
 \frac{\rho_{n+1}}{\rho_n}
 =\frac{(n+1/2)(n+5/2)}{(n+1)(n+2)}<1.
\]
Lemma \ref{lem:sign} therefore proves that all three quotients have
constant coefficient one and strictly negative coefficients
at every positive degree.

\paragraph{The upper remainder has positive coefficients.}
First, Lemma \ref{lem:sign} gives \(1/F_1=1-C\), where every
coefficient of \(C\), starting at degree one, is positive.
For \(0<x<1\), \(C(x)<1\). Thus
\[
 L(x)=-\log(1-C(x))=\sum_{j\geq1}\frac{C(x)^j}{j}
\]
has strictly positive coefficients at all positive degrees.
The same argument applies to \(A=\log F_0\).

Next set
\[
 P_0(x)=\frac{(1-x)^{-1/4}}{F_0(x)},\qquad
 P_1(x)=\frac{(1-x)^{-1/3}}{F_1(x)}.
\]
By \eqref{eq:derivatives},
\[
 (\log P_0)'=\frac{1-G_0/F_0}{4(1-x)},\qquad
 (\log P_1)'=\frac{1-G_1/F_1}{3(1-x)}.
\]
Both right sides have nonnegative coefficients. Integration from
zero and exponentiation show that \(P_0,P_1\) have nonnegative
coefficients and constant coefficient one. Consequently so does
\begin{equation}\label{eq:positivefactor}
 \frac1{(1-x)F_0F_1}=(1-x)^{-5/12}P_0P_1.
\end{equation}

Introduce the source's auxiliary function
\[
 J(x)=F_0(x)-F_1(x)\frac{F_0'(x)}{F_1'(x)}
     =F_0(x)-\frac34F_1(x)\frac{G_0(x)}{G_1(x)}.
\]
Differentiation cancels the first-order terms:
\[
 J'(x)=-\frac34F_1(x)\left(\frac{G_0(x)}{G_1(x)}\right)'.
\]
The preceding quotient-sign result proves that every coefficient
of \(J'\) is strictly positive, and \(J(0)=1/4\).
Since \(a_n<a_{n-1}\) for \(n\geq1\), the series \((1-x)F_0\)
has constant coefficient one and negative coefficients thereafter.
It follows that
\[
 T(x)=4J(x)-(1-x)F_0(x)
\]
has zero constant coefficient and strictly positive coefficients
at every positive degree.

The following identity now closes the proof for \(h_{11}\):
\begin{equation}\label{eq:key}
 4H'(x)=L(x)+\frac{G_1(x)T(x)}{3(1-x)F_0(x)F_1(x)}.
\end{equation}
It follows directly by differentiating \(H\) and using
\eqref{eq:derivatives}. Every factor on the right has nonnegative
coefficients, by \eqref{eq:positivefactor}; \(L\) supplies a strictly
positive coefficient at every positive degree.
The constant coefficient is zero. As \(H(0)=0\), integration proves
\[
 H(x)=\sum_{n\geq2}d_nx^n,\qquad d_n>0.
\]
Dividing by \(x^2\) proves precisely the required coefficient
positivity of \(h_{11}\).

\paragraph{The other five functions.}
Yang--Zhao's Theorem 1.1 and its coefficient conclusion
(6.11)--(6.15) \cite{YZ} establish
\begin{equation}\label{eq:known}
 Q(x)=\sum_{n\geq0}q_nx^n,\quad
 q_n>0,\quad q_0=\frac34,\quad q_1=\frac1{320},\quad
 q_2=\frac{517}{201600},\qquad 0\leq x<1.
\end{equation}
We use this established result with attribution.
It immediately proves the positive coefficients of \(G\) and
\(f=(Q-q_0)/x\). The previous logarithm argument proves those
of \(f_7=L/x\) and \(h_{13}=A/x\). Finally,
\[
 h_{12}(r)
 =\frac{L(x)}x\,
   \frac{Q(x)-q_0-q_1x}{x^2}
\]
is a product of two series having strictly positive coefficients.
This proves every coefficient assertion.

All series converge on the unit disk. This follows from the
coefficient bounds in Lemma \ref{lem:sign}, the displayed convergent
products and logarithms, and \eqref{eq:known}. Termwise
differentiation is therefore valid locally on \((0,1)\).
For any derivative order in \(r\), infinitely many even monomials
of sufficiently high degree give positive contributions.
This proves strict absolute monotonicity, strict increase,
and strict convexity as stated.

\paragraph{Endpoint values and ranges.}
The first two logarithm coefficients are
\[
 A(x)=\frac x4+\frac{7x^2}{64}+O(x^3),\qquad
 L(x)=\frac x3+\frac{13x^2}{90}+O(x^3).
\]
Thus \(H(x)=79x^2/960+O(x^3)\); also \eqref{eq:known} gives
\[
 h_{12}(0+)=\frac13q_2=\frac{517}{604800}.
\]
As \(x\uparrow1\), the standard elementary/elliptic asymptotics give
\[
 F_0(x)\sim\frac1\pi\log\frac1{1-x},\qquad
 F_1(x)\sim\frac12\log\frac1{1-x}.
\]
Consequently \(L-A\to\log(\pi/2)\), \(L\to\infty\), and
\((1-x)L\to0\). Therefore
\[
 H=L-A-\frac{1-x}{4}L\longrightarrow\log(\pi/2),
\]
while
\[
 A-(3/4+x/320)L
 =\left(1-\frac34-\frac{x}{320}\right)L+(A-L)
 \longrightarrow+\infty.
\]
Strict increase and continuity give both exact open ranges.
Theorem \ref{thm:main}, and hence the entire conjecture, follows.

\paragraph{Task 2: Extend the upper-remainder argument.}
Determine the hypergeometric parameter pairs for which an analogue of
\(H(x)=((3+x)/4)\log F_1(x)-\log F_0(x)\) has positive coefficients
after removing its initial zeros.

\paragraph{Task 3: Optimize bounds from truncated remainders.}
Use the positive coefficients of \(h_{11}\) and \(h_{12}\), together
with their endpoint behavior, to derive sharp two-sided bounds for
\(\mathcal K(r)\) at any prescribed truncation degree.

\paragraph{Task 4: Describe the upper-remainder coefficients.}
Determine asymptotic formulas and any Hausdorff-moment or
log-convexity properties of the sequence \((d_{n+2})_{n\geq0}\).

\paragraph{Task 5: Find a direct proof for the quotient coefficients.}
Seek an elementary coefficient argument for \eqref{eq:known}
that integrates with Lemma \ref{lem:sign} and extends to nearby
hypergeometric parameter families.

\begin{thebibliography}{9}
\bibitem{QBMJ}
S.-L.\ Qiu, Q.\ Bao, X.-Y.\ Ma, and H.-B.\ Jiang,
On a Conjecture Concerning the Approximates of Complete Elliptic
Integral of the First Kind by Inverse Hyperbolic Tangent,
\href{https://arxiv.org/abs/2103.04072}{arXiv:2103.04072v1} (2021),
Conjecture 5.1, p.~21; Lemma 2.2 and equation (4.2).

\bibitem{YZ}
Z.-H.\ Yang and T.-H.\ Zhao,
On the absolute monotonicity related to Alzer--Qiu's conjecture,
\emph{Integral Transforms and Special Functions}, published online
13 April 2026,
\href{https://doi.org/10.1080/10652469.2026.2657542}
{doi:10.1080/10652469.2026.2657542},
Theorem 1.1 and equations (6.11)--(6.15).
\end{thebibliography}
\end{document}
